\section{结论、局限性与分工}

\subsection{结论与局限性}

在统一预处理和固定数据划分下，多模型比较结果表明，
SVM 在测试集上的 Recall 和 F1-score 分别达到
0.7778 和 0.6829，均为基础候选模型中的最高值；
Random Forest 的 ROC AUC 达到 0.8275，
体现出较好的概率排序能力。
XGBoost 的 Accuracy 为 0.7468，与 SVM 持平，
但 Recall 和 F1-score 低于 SVM。
Decision Tree 和 Gradient Boosting 在默认参数下表现相对较弱。

上述结果说明，糖尿病风险预测不能只依据 Accuracy 选择模型。
由于漏判高风险样本可能使潜在阳性个体错过进一步检查，
模型选择应综合考虑 Recall、F1-score、ROC AUC 和分类阈值。
本阶段的比较结果为后续调参和最终模型选择提供了依据：
SVM 和 Random Forest 可作为重点优化对象，
XGBoost 和 LightGBM 则可作为备选增强模型。

本文仍存在若干局限。
首先，Pima 数据集样本量有限，
且样本人群限定为至少 21 岁的 Pima Indian 女性，
模型结论不宜直接推广到其他人群。
其次，数据仅包含 8 个结构化生理指标，
缺少饮食、运动、用药史等更完整的临床信息。
最后，项目模型只能用于课程实验和风险提示，
不能作为医学诊断工具。

后续工作可围绕候选模型开展交叉验证和超参数搜索，
并结合验证集调整分类阈值，
在控制漏判风险的同时平衡 Precision 和 Recall。
此外，还可引入 SHAP 或置换重要性解释非线性模型，
收集更丰富的临床变量，
并继续完善 Demo 部署和用户测试。

\subsection{团队分工}

\begin{table}[t]
  \caption{团队成员分工}
  \label{tab:team_contribution}
  \centering
  \scriptsize
  \begin{tabular}{p{0.20\linewidth}p{0.63\linewidth}}
    \toprule
    成员 & 主要负责内容 \\
    \midrule
    柯文丽 & 数据理解、数据预处理、数据划分与基准模型 \\
    王佳妮 & 交互式问诊系统设计、模型推理流程封装、系统集成\\
    秦琪 & 多模型训练、模型性能对比与结果整理 \\
    何晨东 & 参数调优、最终模型评估与结果可视化分析 \\
    \bottomrule
  \end{tabular}
\end{table}

本项目采用模块化分工方式完成。
柯文丽主要负责数据集字段理解、缺失值与异常值处理、特征标准化以及训练集和测试集划分，并建立基准模型；
秦琪负责训练多种机器学习模型，并从 Accuracy、Precision、Recall、F1 和 AUC 等指标对不同模型进行比较；
何晨东负责对候选模型进行参数调优，完成最终模型评估，并绘制混淆矩阵、ROC 曲线和特征重要性等结果图；
王佳妮负责系统集成与 Demo 展示部分：
设计 Gradio 交互式问诊界面，
将用户输入、预处理对象和模型推理流程封装为统一预测接口，
实现单样本风险概率输出、低/中/高风险等级划分和输入因素解释；
同时整理系统运行方式、模型接入约定和系统演示相关报告内容，
保证训练阶段得到的模型能够以可交互形式展示。
各成员同时参与对应模块的报告与展示内容整理。
